\centerline{\bf \Huge Степень точки}
\centerline{\bf \Huge Level: Hard. Решения}

\problem{1.1}
В окружности $\Omega$ проведена хорда $BC$. В один из отсекаемых ею сегментов вписаны две окружности $\omega_1$ и $\omega_2$. Докажите, что середина дуги $BC$ окружности $\Omega$, не касающейся окружностей $\omega_1$ и $\omega_2$, лежит на радикальной оси окружностей $\omega_1$ и $\omega_2$. \\

\problem{1.2}
Окружности $\Omega$ и $\omega$ касаются друг друга внутренним образом в точке $A$. Проведем в большей окружности $\Omega$ хорду $BC$, касающуюся $\omega$ в точке $K$ (хорда $BC$ не является диаметром $\omega$). Точка $M$ --- середина отрезка $AK$. Докажите, что окружность $(BMC)$ проходит через центр $\omega$. 

\textbf{Решение (1.1)} Докажем, что степень середины $W$ дуги $BC$ окружности $\Omega$ относительно любой окружности $\omega$, вписанной в сегмент, отсекаемой хордой $BC$, постоянная и равна $BW^2$. Действительно, рассмотрим произвольную окружность $\omega$, касающуюся окружности $\Omega$ в точке $A$ и хорды $BC$ в точке $K$. По лемме Архимеда точки $A$, $K$ и $W$ лежат на одной прямой, так что степень точки $W$ относительно окружности $\omega$ равна $QK\cdot WA$. С другой стороны, треугольники $WAB$ и $WBK$ подобны по двум углам (угол $\angle W$ у них общий, а углы $\angle WAB$ и $\angle WBK$ смотрят на равные дуги и потому тоже равны). Отсюда следует, что $WB:WK=WA:WB$ и $WK\cdot WA=WB^2$, что и требовалось доказать.

\textbf{Решение (1.2)} Проведем общую внешнюю касательную к окружностям $\Omega$ и $\omega$ в точке $A$. пусть она пересекает прямую $BC$ в точке $T$. Заметим, что точки $T$, $M$ и $O$ лежат на одной прямой, поэтому достаточно проверить справедливость равенства $TM\cdot TO=TB\cdot TC$. Но обе части этого равенства равны величине $TA^2$. Действительно, $TM\cdot TO=TA^2$ по стандартному подобию в прямоугольном треугольнике $TAO$, а $TB\cdot TC=TA^2$ по теореме о квадрате касательной. Таким образом, $TM\cdot TO=TB\cdot TC$ и точки $B$, $C$, $M$ и $O$ лежат на одной окружности, что и требовалось доказать.

 \problem{2}
 В треугольнике $ABC$ проведены чевианы $AA_1$, $BB_1$ и $CC_1$, пересекающиеся в одной точке. Прямая $B_1C_1$ пересекает окружность $(ABC)$ в точках $P$ и $Q$. Докажите, что окружность $(A_1PQ)$ проходит через середину отрезка $BC$.

\textbf{Решение.} Продлим прямую $B_1C_1$ до пересечения с прямой $BC$ в точке $T$ и обозначим середину отрезка $BC$ через $M$. Нам необходимо доказать равенство $TP\cdot TQ=TA_1\cdot TM$. Пусть $TB=x$, $BA_1=y$ и $BC=a$. Тогда по теоремам Чевы и менелая получаем следующие соотношения: $$\frac{BC_1}{C_1A}\cdot \frac{AB_1}{B_1C}\cdot\frac{CA_1}{A_1B}=1, \quad \frac{BC_1}{C_1A}\cdot\frac{AB_1}{B_1C}\cdot\frac{CT}{TB}=1.$$ Поделив эти равенства друг на друга, мы получим следующиее соотношение между длинами $x$, $y$, $a$: $$\frac{a-y}{y}=\frac{a+x}{x}\quad\Longrightarrow\quad ax=2xy+ay.$$ С другой стороны, соотношение $TP\cdot TQ=TA_1\cdot TM$ переписывается в виде $$x(a+x)=(x+y)(x+a/2)\quad\Longleftrightarrow\quad ax=2xy+ay.$$ Таким образом, из теорем Чевы и Менелая следует соотношение $TP\cdot TQ=TA_1\cdot TM$. Отсюда в свою очередь мы получаем, что точки $P$, $A$, $A_1$ и $M$ лежат на одной окружности, что и требовалось доказать.

 \problem{3}
 В треугольнике $ABC$ отмечена точка пересечения высот $H$. Окружность с центром в середине стороны $AB$ и проходящая через $H$, пересекает сторону $AB$ в точках $C_1$ и $C_2$. Аналогично определяются точки $A_1$, $A_2$ и $B_1$, $B_2$. Докажите, что точки $A_1$, $A_2$, $B_1$, $B_2$, $C_1$, $C_2$ лежат на одной окружности. 

\textbf{Решение 1} Достаточно доказать, что четверки точек $(A_1, A_2, B_1,B_2)$, $(B_1, B_2, C_1, C_2)$ и $(C_1, C_2, A_1,A_2)$ лежат на одной окружности. Действительно, если эти три окружности не совпадают, то их радикальные оси $A_1A_2$, $B_1B_2$ и $C_1C_2$ должны пересекаться в одной точке, что невозможно, т.к. они содержат стороны трегольника $ABC$.

Докажем, что точки $B_1$, $B_2$, $C_1$ и $C_2$ лежат на одной окружности. Для этого докажем, что $AB_1\cdot AB_2=AC_1\cdot AC_2$. Это условие равносильно тому, что точка $A$ лежит на радикальной оси окружностей $\omega_b$ и $\omega_c$ из условия (их центры находятся в серединах отрезков $AC$ и $AB$, а сами они проходят через точку $H$). Иначе говоря, нужно доказать, что прямая $AH$ является радикальной осью окружностей $\omega_b$ и $\omega_c$. Но это очевидно, т.к. прямая $AH$ проходит через точку $H$ пересечения этих окружностей и перпендикулярна их линии центров (которая есть просто средняя линия, параллельная стороне $BC$). 

\textbf{Решение 2} Как и в предыдущем решении, будем доказывать, что точки $B_1$, $B_2$, $C_1$ и $C_2$ лежат на одной окружности и будем проверять равенство $AB_1\cdot AB_2=AC_1\cdot AC_2$. Положим $AC=2b$, $AB=2c$, отметим также середины $M_b$ и $M_c$ сторон $AC$ и $AB$ и длины $M_bH=x$ и $M_cH=y$. Тогда равенство $AB_1\cdot AB_2=AC_1\cdot AC_2$ переписывается следующим образом: $b^2-x^2=c^2-y^2$ или, что то же самое, $b^2-c^2=x^2-y^2$. 

Положим $BH=u$, $CH=v$, $AH=t$ и запишем формулу длины медианы в треугольниках $ABH$ и $ACH$: $$y^2=\frac{2(u^2+t^2)-4c^2}{4}, \quad x^2=\frac{2(v^2+t^2)-4b^2}{4}.$$ Вычитая из второго равенства первое и учитывая, что $v^2-u^2=4b^2-4c^2$ по теореме Карно, получаем: $$x^2-y^2=\frac{v^2-u^2}{2}-(b^2-c^2)=\frac{4(b^2-c^2)}{2}-(b^2-c^2)=b^2-c^2,$$ что и требовалось.

 \problem{4}
 На стороне $AB$ выпуклого четырёхугольника $ABCD$ взяты точки $K$ и $L$ (точка $K$ лежит между $A$ и $L$), а на стороне $CD$ взяты точки $M$ и $N$ (точка $M$ между $C$ и $N$). Известно, что $AK=KN=ND$ и $LB=BC=CM$. Докажите, что если $BCNK$ --- вписанный четырехугольник, то и $ADML$ тоже вписан.

\textbf{Решение.} Продлим стороны $AB$ и $CD$ до пересечения в точке $P$. Пусть $LB=BC=CM=x$, $AK=KN=ND=y$, $KL=u$, $MN=v$, $PB=a$ и $PC=b$. Запишем подобие треугольников $PBC$ и $PNK$: $$\frac{a}{b+x+v}=\frac{b}{a+x+u}=\frac{x}{y}\quad\Longleftarrow\quad a+x+u=\frac{by}{x}, \quad b+x+v=\frac{ay}{x}.$$ Для того, чтобы четырехугольник $ALMD$ был вписанным, необходмо и достаточно, чтобы $PL\cdot PA=PM\cdot PD$. Запишем это равенство в переменных $x$, $y$, \ldots и подставим в него формулы, полученные выше:
\begin{multline*}
  (a+x)(a+x+u+y) =^? (b+x)(b+x+v+y)\quad\Longleftrightarrow \\ \Longleftrightarrow\quad (a+x)\Big(\frac{by}{x}+y\Big) =^? (b+x)\Big(\frac{ay}{x}+y\Big)  \quad\Longleftrightarrow \\ \Longleftrightarrow\quad \frac{y}{x}\cdot (a+x)(b+x) = ^?\frac{y}{x}\cdot (a+x)(b+x).
\end{multline*}

Последнее равенство верно, а значит, верны и все предыдущие равенства. Таким образом, мы доказали, что верно равенство $PL\cdot PA=PM\cdot PD$, а из него немедленно следует вписанность четырехугольника $ADML$.

 \problem{5}
 В остроугольном треугольнике $ABC$ высоты $AA_1$, $BB_1$ и $CC_1$ пересекаются в точке $H$. Точка $O$ --- центр окружности $(ABC)$. Докажите, что точки $O$, $H$, $A_1$ и точка, симметричная точке $A$ относительно прямой $B_1C_1$, лежат на одной окружности.

\textbf{Решение.} Обозначим точку, симметричную точке $A$ относительно прямой $B_1C_1$, через $A'$. Прежде всего заметим, что точка $A'$ лежит на прямой $AO$, поскольку прямые $B_1C_1$ и $AO$ перпендикулярны. Поэтому задача следует из справедливости равенства $AH\cdot AA_1=AA'\cdot AO$. Положим $AO=R$, $AA_1=h$ и $\angle BAC=\alpha$. Треугольники $AB_1C_1$ и $ACB$ подобны с коэффициентом $\cos\alpha$, поэтому $AA'=2h\cos\alpha$. Далее, если $M$ --- середина стороны $BC$, то $AH=2OM$, а из треугольника $COM$ следует, что $OM=R\cos\alpha$ (напомним, что $\angle COM=\angle COB/2=\alpha$). Таким образом, оба произведения $AH\cdot AA_1$ и $AA'\cdot AO$ равны одной и той же величине $2hR\cos\alpha$. Отсюда и следует наше утверждение.

 \problem{6}
 Окружность $\omega$ касается сторон угла $\angle BAC$ в точках $B$ и $C$. Прямая $\ell$ пересекает отрезки $AB$ и $AC$ в точках $K$ и $L$ соответственно. Окружность $\omega$ пересекает прямую $\ell$ в точках $P$ и $Q$. Точки $S$ и $T$ выбраны на отрезке $BC$ так, что $KS$ параллельно $AC$ и $LT$ параллельно $AB$. Докажите, что точки $P$, $Q$, $S$ и $T$ лежат на одной окружности.

\textbf{Решение.} По теореме о пропорциональных отрезках имеем $XP\cdot XQ=XB\cdot XC$. Значит, для решения задачи нам достаточно проверить справедливость равенства $XB\cdot XC=XS\cdot XT$. Для этого заметим, что треугольники $KBS$ и $LCT$ равнобедренные: $KB=KS$ и $LC=LT$. Далее, запишем подобие пар треугольников $XKS$, $XLC$ и $XKB$, $XLT$: $$\frac{XS}{XC}=\frac{KS}{LC}=\frac{KB}{LC}, \quad \frac{XB}{XT}=\frac{KB}{LT}=\frac{KB}{LC}.$$ Значит, $XS/XC=XB/XT$ и $XB\cdot XC=XS\cdot XT$, что и требовалось доказать.

\problem{7}Для окружности $\omega$ и точки $X$ обозначим через $\deg_\omega X$ степень этой точки относительно окружности $\omega$. Докажите, что для любых прямых $AC$ и $BD$, пересекающихся в точке $P$, и любой точки $X$ имеет место равенство $$\deg_{(PAB)}(X)+\deg_{(PCD)}(X)=\deg_{(PAD)}(X)+\deg_{(PBC)}(X).$$

\textbf{Решение.} Заметим, что функция $$f(X)=Pow_{(PAB)}(X)+Pow_{(PCD)}(X)-Pow_{(PAD)}(X)-Pow_{(PBC)}(X)$$ является линейной и обращается в 0 в точках $A$, $B$ и $C$, не лежащих на одной прямой. Значит, эта функция тождественно равна 0, что и требовалось доказать.

\problem{8} Пусть $A$, $V_1$, $V_2$, $B$, $U_2$, $U_1$ --- точки на окружности $\Omega$, выбранные в таком порядке, что $BU_2>AU_1>BV_2>AV_1$. Точка $X$ --- произвольная точка на дуге $V_1V_2$, не содержащей точки $A$ и $B$. Пусть $XA\cap U_1V_1=C$, $XB\cap U_2V_2=D$. Докажите, что существует такая точка $K$, не зависящая от точки $X$, что степень этой точки относительно окружности $(XCD)$ постоянна при любом выборе точки $X$.

\textbf{Решение.} Забудем про условие принадлежности точки $X$ дуге $V_1V_2$ и будем рассматривать точку $X$ на высех окружности $\Omega$. Тогда выберем точку $E$ на $\Omega$ так, чтобы прямые $AE$ и $U_1V_1$ были бы параллельны. В таком случае точка $C$ станет бесконечно удаленной, а окружность $(XCD)$ при $X=E$ превратится в прямую $BE$. Аналогично, выбрав на окружности $\Omega$ такую точку $F$, что прямые $BF$ и $U_2V_2$ параллельны, мы получим, что окружность $(XCD)$ при $X=F$ превратится в прямую $AF$. Значит, на роль точки $K$ может подойти только точка пересечения прямых $AF$ и $BE$: ее степень относительно вырожденных <<окружностей>> $AF$ и $BE$ равна 0. Докажем, что ее степень относительно любой окружности $(XCD)$ будет одинакова.

Для этого используем п. {\bf ( a )} и заметим, что $$\deg_{(XCD)}(K)=\deg_{(XAD)}(K)+\deg_{(XBC)}(K)-\deg_{(XAB)}(K).$$ Окружность $(XAB)$ совпадает с окружностью $\Omega$, так что последнее слагаемое является константой. Докажем, что первые два слагаемых также постоянны. Для этого отметим точку $P$ пересечения прямых $U_1V_1$ и $BE$, а также точку $Q$ пересечения прямых $U_2V_2$ и $AF$. Заметим, что четырехугольник $BPCX$ --- вписанный: прямые $СP$ и $AE$ параллельны, а прямые $AE$ и $BX$ антипараллельны в угле $\angle (AV_1, BE)$, поэтому прямые $CP$ и $BX$ также антипараллельны в угле $\angle (AV_1, BE)=\angle (CX, BP)$. Аналогично, четырехугольник $AQDX$ является вписанным. Но тогда $\deg_{(XAD)}(K)=KA\cdot KQ=\mathrm{const}$ и $\deg_{(XBC)}(K)=KB\cdot KP=\mathrm{const}$, а значит, и $\deg_{(XCD)}(K)=\mathrm{const}$ что и требовалось.

